% ===== Topic 2, part b: problems 16-30 (VI) =====
% helper: 2x2 grid with one letter; #1,#2 = centre of the cell, #3 = letter, #4 = node transform
\providecommand{\tbGrid}[4]{\begin{tikzpicture}[baseline=0.2cm,x=0.5cm,y=0.5cm]
  \draw[Ink,line width=0.7pt] (0,0) grid[step=1] (2,2);
  \node[font=\sffamily\bfseries,text=tmBlue!80!black,inner sep=0pt,#4] at (#1,#2) {#3};\end{tikzpicture}}

% ---------------------------------------------------------------- 16
\begin{probbox}{16}{\relax}
Một mảnh vườn hình chữ nhật có chiều dài hơn \emph{hai lần} chiều rộng là $6$ m.
Chu vi của mảnh vườn là $120$ m. Hỏi chiều rộng của mảnh vườn bằng bao nhiêu mét?
\choices{$8$ m}{$18$ m}{$27$ m}{$38$ m}{$22$ m}
\end{probbox}
\probsol{
\textbf{Phân tích đề.} Ta biết hai điều: (1) chiều dài $=$ hai lần chiều rộng, cộng thêm $6$ m;
(2) chu vi $=120$ m. Cần tìm chiều rộng. Ta gọi ẩn và lập phương trình.

\step{1}{Gọi ẩn và biểu diễn chiều dài theo ẩn}
Gọi chiều rộng của mảnh vườn là $x$ (m), điều kiện $x>0$.
``Hai lần chiều rộng'' là $2x$; chiều dài hơn hai lần chiều rộng $6$ m nên chiều dài là $2x+6$ (m).
\begin{figC}
\draw[Ink,thick,fill=tmGreen!10] (0,0) rectangle (4.2,1.8);
\node[below] at (2.1,0){$2x+6$};
\node[right] at (4.2,0.9){$x$};
\end{figC}

\begin{recallbox}Chu vi hình chữ nhật $=(\text{dài}+\text{rộng})\cdot 2$.\end{recallbox}

\step{2}{Lập phương trình}
Chu vi bằng $120$ m nên
\[ \big((2x+6)+x\big)\cdot 2=120. \]

\step{3}{Giải phương trình}
Chia hai vế cho $2$: $(2x+6)+x=60$.
Thu gọn vế trái: $3x+6=60$.
Chuyển $6$ sang vế phải: $3x=60-6=54$.
Chia hai vế cho $3$: $x=54:3=18$ (thỏa mãn $x>0$).

Vậy chiều rộng là $18$ m, chiều dài là $2\cdot 18+6=42$ (m).

\begin{warnbox}Nhiều bạn quên rằng chu vi gồm \emph{hai} chiều dài và \emph{hai} chiều rộng, nên viết
$3x+6=120$ và được $x=38$ -- đó chính là phương án (D) ``bẫy''.\end{warnbox}

\verify{Chiều dài $42$ m, chiều rộng $18$ m: $42=2\cdot 18+6$ (đúng) và chu vi
$(42+18)\cdot 2=60\cdot 2=120$ m (đúng).}
Chiều rộng $18$ m ứng với phương án \textbf{(B)}.
\ans{(B) $18$ m}
}

% ---------------------------------------------------------------- 17
\begin{probbox}{17}{\relax}
Chữ P được viết trong một lưới ô vuông $2\times 2$ như hình bên: \tbGrid{1.5}{1.5}{P}{}\,.
Người ta thực hiện liên tiếp một số phép quay quanh tâm của lưới và một số phép lật (đối xứng)
qua hai đường thẳng đi qua tâm lưới (đường nằm ngang và đường thẳng đứng), và nhận được kết quả:
\tbGrid{0.5}{0.5}{P}{cm={0,-1,-1,0,(0,0)}}\,.

Nếu thực hiện \emph{đúng dãy phép biến hình đó} với lưới \tbGrid{0.5}{1.5}{A}{} thì ta được hình nào?
\par\smallskip\noindent
\begin{tabularx}{\linewidth}{@{}XXXXX@{}}
(A) \tbGrid{0.5}{1.5}{A}{} & (B) \tbGrid{0.5}{1.5}{A}{rotate=90} & (C) \tbGrid{1.5}{0.5}{A}{rotate=-90} &
(D) \tbGrid{1.5}{0.5}{A}{} & (E) \tbGrid{1.5}{0.5}{A}{rotate=180}
\end{tabularx}
\end{probbox}
\probsol{
\textbf{Phân tích đề.} Ta không biết dãy phép biến hình cụ thể gồm những bước nào, nhưng biết
\emph{kết quả cuối cùng} của nó trên chữ P. Ý tưởng: tưởng tượng lưới là một \emph{tấm thẻ} vuông
bằng giấy kính trong, có vẽ chữ trên đó. Mỗi phép quay hay lật chỉ là xoay hoặc lật tấm thẻ rồi đặt lại
đúng chỗ cũ. Cả dãy phép biến hình cũng chỉ là \emph{một cách đặt lại tấm thẻ}. Ta đọc cách đặt đó từ
hình chữ P, rồi áp dụng cho chữ A.

\step{1}{Quan sát chữ P đã đi đâu}
Chữ P ở ô \emph{trên -- phải} đã chuyển sang ô \emph{dưới -- trái}. Đồng thời, nét thẳng (nét ``thân'')
của chữ P, vốn thẳng đứng, nay nằm ngang ở phía trên; phần ``bụng'' tròn vốn ở bên phải thân,
nay nằm ở phía dưới.

\step{2}{Tìm xem tấm thẻ đã được đặt lại như thế nào}
Có hai cách đặt lại tấm thẻ đưa ô trên -- phải về ô dưới -- trái:
\begin{itemize}[leftmargin=1.4em,itemsep=1pt]
\item quay $180^\circ$: khi đó chữ P sẽ bị lộn ngược (thân vẫn \emph{thẳng đứng}, trông giống chữ ``d'');
\item lật qua đường chéo nối góc trên -- trái với góc dưới -- phải (đường nét đứt trong hình dưới):
khi đó thân thẳng đứng của P sẽ thành \emph{nằm ngang}.
\end{itemize}
Hình đề bài cho thân nằm ngang, nên dãy phép biến hình có tác dụng như
\textbf{phép lật qua đường chéo trên-trái -- dưới-phải}.
\begin{figC}\begin{scope}[x=0.8cm,y=0.8cm]
\draw[Ink,line width=0.8pt] (0,0) grid[step=1] (2,2);
\draw[tmOrange,dashed,line width=1pt] (-0.3,2.3)--(2.3,-0.3);
\node[font=\sffamily\bfseries\large,text=tmBlue] at (1.5,1.5){P};
\node[font=\sffamily\bfseries\large,text=tmGreen!80!black,cm={0,-1,-1,0,(0,0)}] at (0.5,0.5){P};
\draw[->,tmPurple,thick] (1.55,1.15) to[bend left=35] (0.85,0.45);
\node[font=\footnotesize,text=tmOrange!80!black,anchor=west] at (2.3,-0.3){trục lật};
\begin{scope}[xshift=4.2cm]
\draw[Ink,line width=0.8pt] (0,0) grid[step=1] (2,2);
\draw[tmOrange,dashed,line width=1pt] (-0.3,2.3)--(2.3,-0.3);
\node[font=\sffamily\bfseries\large,text=tmBlue] at (0.5,1.5){A};
\node[font=\footnotesize,text=Slate] at (1,-0.75){A nằm ngay trên trục};
\end{scope}
\end{scope}\end{figC}
\begin{recallbox}Một dãy phép quay quanh tâm và phép lật qua đường ngang, đường dọc có thể cho kết quả giống
một phép lật qua đường chéo. Chẳng hạn: \emph{quay $90^\circ$ theo chiều kim đồng hồ rồi lật qua đường
thẳng đứng}. Thử với P: quay $90^\circ$ đưa P từ ô trên -- phải xuống ô dưới -- phải, lật qua đường
dọc đưa P sang ô dưới -- trái, và thân P nằm ngang -- đúng như hình đề bài.\end{recallbox}

\step{3}{Áp dụng cho chữ A}
Ô trên -- trái nằm \emph{ngay trên} trục lật (đường chéo đi qua nó), nên chữ A \emph{vẫn ở ô trên -- trái}.
Phép lật qua đường chéo này biến hướng ``lên trên'' thành hướng ``sang trái''. Đỉnh nhọn của chữ A vốn
chỉ lên trên, nay chỉ sang trái. Chữ A đối xứng trái -- phải nên không cần lo chuyện bị ``soi gương''.

\begin{warnbox}Đừng chỉ nhìn vị trí ô! Chữ P đi từ ô trên -- phải sang ô dưới -- trái, nhưng chữ A
\emph{không} vì thế mà cũng ``đi chéo xuống''. Phải tìm phép biến hình tác động lên \emph{cả tấm thẻ},
rồi mới áp dụng.\end{warnbox}

\verify{Dùng dãy cụ thể ``quay $90^\circ$ theo chiều kim đồng hồ rồi lật qua đường dọc'': quay đưa
A từ ô trên -- trái sang ô trên -- phải, đỉnh A chỉ sang phải; lật qua đường dọc đưa A về ô trên -- trái,
đỉnh A chỉ sang trái. Đúng là hình (B).}
\ans{(B) chữ A ở ô trên -- trái, đỉnh chỉ sang trái}
}

% ---------------------------------------------------------------- 18
\begin{probbox}{18}{\relax}
Trong hình vuông cạnh $8$ cm, người ta vẽ hai đường chéo. Diện tích phần tô màu bằng bao nhiêu?
\begin{figC}\begin{scope}[x=0.38cm,y=0.38cm]
\fill[tmBlue!35] (0,0)--(8,0)--(4,4)--cycle;
\draw[Ink,line width=1pt] (0,0) rectangle (8,8);
\draw[Ink,line width=0.9pt] (0,0)--(8,8) (0,8)--(8,0);
\node[left] at (0,4){$8$ cm};
\end{scope}\end{figC}
\choices{$4$ cm$^2$}{$8$ cm$^2$}{$16$ cm$^2$}{$56$ cm$^2$}{$64$ cm$^2$}
\end{probbox}
\probsol{
\textbf{Phân tích đề.} Hai đường chéo chia hình vuông thành $4$ tam giác. Phần tô màu là một trong
bốn tam giác đó. Cần tìm diện tích của nó.

\step{1}{Tính diện tích hình vuông}
\begin{recallbox}Diện tích hình vuông cạnh $a$ là $a\times a$.\end{recallbox}
$S_{\text{vuông}}=8\times 8=64$ (cm$^2$).

\step{2}{Chứng tỏ bốn tam giác có diện tích bằng nhau}
\begin{proofbox}Hai đường chéo của hình vuông bằng nhau và cắt nhau tại trung điểm $O$ của mỗi đường, nên
$O$ cách đều bốn đỉnh. Quay hình vuông $90^\circ$ quanh $O$ thì mỗi tam giác chồng khít lên tam giác
bên cạnh. Vậy bốn tam giác bằng nhau, mỗi tam giác chiếm $\tfrac14$ hình vuông.\end{proofbox}

\step{3}{Tính diện tích phần tô màu}
$S=64:4=16$ (cm$^2$).

\verify{Tính trực tiếp: tam giác tô màu có đáy là cạnh dưới dài $8$ cm, chiều cao là khoảng cách từ
tâm $O$ xuống cạnh đáy, bằng nửa cạnh, tức $4$ cm. Diện tích $=\dfrac{8\times 4}{2}=16$ cm$^2$. Khớp.}
Phương án \textbf{(C)}.
\ans{(C) $16$ cm$^2$}
}

% ---------------------------------------------------------------- 19
\begin{probbox}{19}{\relax}
Lưới ô vuông dưới đây gồm các ô vuông nhỏ bằng nhau. Phần tô màu chiếm bao nhiêu phần trăm diện tích
cả lưới?
\begin{figC}\begin{scope}[x=0.36cm,y=0.36cm]
\foreach \r in {0,...,7}{\foreach \c in {0,...,7}{\fill[tmBlue!40] (\c,\r) rectangle (\c+1,\r+1);}}
\foreach \c/\r in {2/7,5/7,2/6,5/6,0/5,1/5,6/5,7/5,0/2,1/2,6/2,7/2,2/1,5/1,2/0,5/0}{\fill[white] (\c,\r) rectangle (\c+1,\r+1);}
\draw[Ink,line width=0.7pt] (0,0) grid[step=1] (8,8);
\end{scope}\end{figC}
\choices{$84\%$}{$80\%$}{$90\%$}{$75\%$}{$66\%$}
\end{probbox}
\probsol{
\textbf{Phân tích đề.} Muốn tính phần trăm, ta cần biết: có tất cả bao nhiêu ô, và có bao nhiêu ô tô màu.
Đếm ô \emph{trắng} sẽ nhanh hơn vì chúng ít hơn.

\step{1}{Đếm tổng số ô}
Lưới có $8$ hàng, mỗi hàng $8$ ô: $8\times 8=64$ ô.

\step{2}{Đếm số ô trắng}
Các ô trắng xếp thành $8$ ``cặp'' hai ô liền nhau, nằm đối xứng ở bốn phía:
\begin{itemize}[leftmargin=1.4em,itemsep=1pt]
\item hai hàng trên cùng: cột thứ $3$ và cột thứ $6$ có ô trắng ở cả hai hàng: $4$ ô;
\item hai hàng dưới cùng: tương tự, $4$ ô;
\item hàng thứ $3$ (từ trên xuống): $2$ ô trắng sát mép trái và $2$ ô sát mép phải: $4$ ô;
\item hàng thứ $6$: tương tự, $4$ ô.
\end{itemize}
Tổng số ô trắng: $4+4+4+4=16$ ô.

\step{3}{Tính số ô tô màu và đổi ra phần trăm}
Số ô tô màu: $64-16=48$ ô.
\begin{recallbox}Tỉ số phần trăm của $a$ so với $b$ là $\dfrac{a}{b}\times 100\%$.\end{recallbox}
\[\frac{48}{64}=\frac{3}{4}=\frac{75}{100}=75\%.\]

\begin{tipbox}Có thể không cần đếm hết: ô trắng chiếm $\dfrac{16}{64}=\dfrac14=25\%$, nên ô tô màu chiếm
$100\%-25\%=75\%$.\end{tipbox}

\verify{$75\%$ của $64$ là $64\times\dfrac34=48$ ô, và $48+16=64$. Đúng.}
Phương án \textbf{(D)}.
\ans{(D) $75\%$}
}

% ---------------------------------------------------------------- 20
\begin{probbox}{20}{\relax}
Cho $JKLM$ là hình vuông và $PQRS$ là hình chữ nhật, đặt chồng lên nhau như hình vẽ. Biết $JK$ song song
với $PQ$, $JK=8$ và $PS=2$. Tính tổng diện tích các phần tô màu.
\begin{figC}\begin{scope}[x=0.3cm,y=0.3cm]
\fill[tmBlue!35] (5,-3) rectangle (13,5);
\draw[Ink,line width=0.9pt] (5,-3) rectangle (13,5);
\fill[white] (0,0) rectangle (18,2);
\draw[Ink,line width=0.9pt] (0,0) rectangle (18,2);
\node[above left] at (0,2){$P$};\node[above right] at (18,2){$Q$};
\node[below right] at (18,0){$R$};\node[below left] at (0,0){$S$};
\node[above] at (5,5){$J$};\node[above] at (13,5){$K$};
\node[below] at (13,-3){$L$};\node[below] at (5,-3){$M$};
\end{scope}\end{figC}
\choices{$32$}{$16$}{$56$}{$48$}{$62$}
\end{probbox}
\probsol{
\textbf{Phân tích đề.} Phần tô màu là phần của hình vuông \emph{không} bị dải hình chữ nhật che.
Vậy: diện tích tô màu $=$ diện tích hình vuông $-$ diện tích phần chung của hai hình.

\step{1}{Tính diện tích hình vuông $JKLM$}
\begin{recallbox}Diện tích hình vuông $=$ cạnh $\times$ cạnh; diện tích hình chữ nhật $=$ dài $\times$ rộng.\end{recallbox}
$S_{JKLM}=8\times 8=64$.

\step{2}{Xác định phần chung}
Vì $JK\parallel PQ$ nên các cạnh của dải $PQRS$ song song với cạnh $JK$ của hình vuông, tức là dải
nằm ``ngang'' cắt qua hình vuông. Dải dài hơn hình vuông nên cắt ngang trọn bề rộng hình vuông.
Phần chung là một hình chữ nhật có:
\begin{itemize}[leftmargin=1.4em,itemsep=1pt]
\item chiều dài bằng cạnh hình vuông: $8$;
\item chiều rộng bằng bề rộng của dải: $PS=2$.
\end{itemize}
Diện tích phần chung: $8\times 2=16$.

\step{3}{Tính diện tích tô màu}
$64-16=48$.

\begin{warnbox}Không cần biết dải nằm cao hay thấp, và hai mảnh tô màu to nhỏ ra sao: tổng của chúng
luôn là $64-16=48$. Nhiều bạn chỉ tính $8\times 2=16$ (phương án B) -- đó là phần \emph{bị che}, không phải
phần tô màu.\end{warnbox}

\verify{Giả sử dải cách cạnh trên $JK$ một khoảng $3$; mảnh trên có diện tích $8\times 3=24$, mảnh dưới có
chiều cao $8-3-2=3$, diện tích $24$. Tổng $48$. Đúng.}
Phương án \textbf{(D)}.
\ans{(D) $48$}
}

% ---------------------------------------------------------------- 21
\begin{probbox}{21}{\relax}
Ta dựng được tam giác có ba cạnh dài $4$, $5$, $8$, nhưng \emph{không} dựng được tam giác có ba cạnh
$4$, $5$, $10$. Chỉ dùng các độ dài $2$, $3$, $5$, $7$, $11$ (mỗi độ dài có thể được dùng nhiều lần), có thể dựng
được bao nhiêu tam giác khác nhau có \emph{đúng hai cạnh bằng nhau}?
\choices{$8$}{$5$}{$20$}{$10$}{$14$}
\end{probbox}
\probsol{
\textbf{Phân tích đề.} Tam giác có đúng hai cạnh bằng nhau là tam giác cân nhưng không đều: hai cạnh bên
bằng $a$, cạnh đáy bằng $b$ với $b\neq a$. Ta phải đếm các cặp $(a,b)$ lấy từ $\{2,3,5,7,11\}$ sao cho ba
đoạn $a,a,b$ thật sự dựng được tam giác.

\begin{recallbox}\textbf{Bất đẳng thức tam giác.} Ba đoạn thẳng là ba cạnh của một tam giác khi và chỉ khi
tổng hai đoạn bất kì lớn hơn đoạn còn lại. Thực hành: chỉ cần kiểm tra
\emph{tổng hai đoạn ngắn hơn lớn hơn đoạn dài nhất}. Ví dụ: $4+5=9>8$ nên dựng được; $4+5=9<10$ nên không
dựng được.\end{recallbox}

\step{1}{Rút gọn điều kiện cho tam giác cân $a,a,b$}
Các tổng cần kiểm tra là $a+a>b$ và $a+b>a$. Điều thứ hai luôn đúng (vì $b>0$). Vậy chỉ cần
\[ 2a>b \quad\text{(hai lần cạnh bên lớn hơn cạnh đáy).}\]

\step{2}{Lập bảng theo cạnh bên $a$}
\begin{center}
\begin{tabular}{|c|c|c|c|}
\hline
Cạnh bên $a$ & $2a$ & Cạnh đáy $b\ne a$ thỏa $b<2a$ & Số tam giác\\ \hline
$2$ & $4$ & $3$ & $1$\\
$3$ & $6$ & $2,\ 5$ & $2$\\
$5$ & $10$ & $2,\ 3,\ 7$ & $3$\\
$7$ & $14$ & $2,\ 3,\ 5,\ 11$ & $4$\\
$11$ & $22$ & $2,\ 3,\ 5,\ 7$ & $4$\\ \hline
\end{tabular}
\end{center}
Chẳng hạn với $a=3$: đáy $b=7$ không được vì $3+3=6<7$; $b=11$ cũng không được.
Với $a=5$: $b=11$ không được vì $5+5=10<11$.

\step{3}{Cộng lại}
$1+2+3+4+4=14$ tam giác.

\begin{warnbox}Tam giác $2,2,5$ hay $3,3,7$ ``trông có vẻ'' được nhưng thật ra hai cạnh bên quá ngắn,
không chạm nhau được. Ngoài ra, tam giác $5,5,5$ có \emph{ba} cạnh bằng nhau nên không được đếm.\end{warnbox}

\verify{Đếm ngược: có $5\times 4=20$ cặp $(a,b)$ với $b\neq a$. Các cặp bị loại (vì $b\ge 2a$) là
$(2;5),(2;7),(2;11),(3;7),(3;11),(5;11)$: $6$ cặp. $20-6=14$. Khớp.}
Phương án \textbf{(E)}.
\ans{(E) $14$}
}

% ---------------------------------------------------------------- 22
\begin{probbox}{22}{\relax}
Hình bình hành $PQRS$ có diện tích bằng $40$. Gọi $T$ và $V$ lần lượt là trung điểm của các cạnh $PS$
và $RS$. Tính diện tích tứ giác $PRVT$.
\begin{figC}\begin{scope}[x=0.55cm,y=0.55cm]
\fill[tmBlue!25] (0,0)--(8,4)--(7,2)--(3,0)--cycle;
\draw[Ink,line width=1pt] (0,0)--(2,4)--(8,4)--(6,0)--cycle;
\draw[Ink,dashed,line width=0.8pt] (0,0)--(8,4) (3,0)--(7,2);
\foreach \x in {1.5,4.5}{\draw[Ink] (\x,-0.15)--(\x,0.15);}
\foreach \x/\y in {6.5/1,7.5/3}{\draw[Ink] (\x-0.15,\y+0.07)--(\x+0.15,\y-0.07);\draw[Ink] (\x-0.15,\y+0.17)--(\x+0.15,\y+0.03);}
\node[below left] at (0,0){$P$};\node[above left] at (2,4){$Q$};
\node[above right] at (8,4){$R$};\node[below right] at (6,0){$S$};
\node[below] at (3,0){$T$};\node[right] at (7,2){$V$};
\end{scope}\end{figC}
\choices{$10$}{$12$}{$15$}{$16$}{$18$}
\end{probbox}
\probsol{
\textbf{Phân tích đề.} Tứ giác $PRVT$ là phần còn lại khi lấy tam giác $PRS$ bỏ đi tam giác nhỏ $TVS$
ở góc $S$. Vậy ta tính $S_{PRS}$ và $S_{TVS}$.

\step{1}{Diện tích tam giác $PRS$}
\begin{recallbox}Đường chéo chia hình bình hành thành hai tam giác bằng nhau, nên mỗi tam giác có diện tích
bằng nửa diện tích hình bình hành.\end{recallbox}
$S_{PRS}=40:2=20$.

\step{2}{Diện tích tam giác $TRS$ (nối thêm $T$ với $R$)}
\begin{recallbox}Hai tam giác có \emph{chung chiều cao} thì tỉ số diện tích bằng tỉ số hai đáy. (Vì diện tích
tam giác $=\tfrac12\times$ đáy $\times$ chiều cao.)\end{recallbox}
Tam giác $TRS$ và tam giác $PRS$ có chung đỉnh $R$, hai đáy $TS$ và $PS$ nằm trên cùng một đường thẳng,
nên có chung chiều cao hạ từ $R$. Vì $T$ là trung điểm $PS$ nên $TS=\tfrac12 PS$. Do đó
$S_{TRS}=\tfrac12\cdot 20=10$.

\step{3}{Diện tích tam giác $TVS$}
Tam giác $TVS$ và tam giác $TRS$ có chung đỉnh $T$, đáy $VS$ và $RS$ cùng nằm trên đường thẳng $RS$.
Vì $V$ là trung điểm $RS$ nên $VS=\tfrac12 RS$, suy ra $S_{TVS}=\tfrac12\cdot 10=5$.

\step{4}{Diện tích tứ giác $PRVT$}
$S_{PRVT}=S_{PRS}-S_{TVS}=20-5=15$.

\begin{tipbox}Tam giác ``nửa -- nửa'' ở góc có diện tích bằng $\tfrac12\times\tfrac12=\tfrac14$ tam giác lớn.
Đây là một kết quả rất hay dùng!\end{tipbox}

\verify{Đặt hình lên lưới: $P(0;0)$, $S(6;0)$, $R(8;4)$, $Q(2;4)$ (diện tích $6\times4=24$, không phải $40$,
nhưng tỉ lệ thì như nhau). Khi đó $T(3;0)$, $V(7;2)$; $S_{PRS}=12$, $S_{TVS}=\tfrac12\cdot 3\cdot 2=3$,
nên $S_{PRVT}=9=\tfrac{3}{8}\cdot 24$. Với hình đề bài: $\tfrac38\cdot 40=15$. Khớp.}
Phương án \textbf{(C)}.
\ans{(C) $15$}
}

% ---------------------------------------------------------------- 23
\begin{probbox}{23}{\stars{1}}
Trên lưới ô vuông $10\times 10$ (mỗi ô có cạnh $1$), cho tam giác $FGH$ như hình vẽ. Điểm $P$ có thể là
một trong $41$ giao điểm của các đường kẻ lưới nằm \emph{bên trong} (không nằm trên cạnh) tam giác $FGH$.
Với mỗi vị trí của $P$ ta được đúng ba tam giác: $FPG$, $GPH$, $HPF$. Trong tất cả $123$ tam giác đó, có
bao nhiêu tam giác có diện tích \emph{đúng bằng một nửa} diện tích tam giác $FGH$?
\begin{figC}\begin{scope}[x=0.36cm,y=0.36cm]
\draw[Line,line width=0.5pt] (0,0) grid[step=1] (10,10);
\draw[Ink,line width=1.3pt] (0,0)--(10,0)--(5,10)--cycle;
\node[above] at (5,10){$F$};\node[below left] at (0,0){$G$};\node[below right] at (10,0){$H$};
\end{scope}\end{figC}
\choices{$5$}{$7$}{$3$}{$11$}{$9$}
\end{probbox}
\probsol{
\textbf{Phân tích đề.} Đặt hệ trục: $G(0;0)$, $H(10;0)$, $F(5;10)$. Tam giác $FGH$ có đáy $GH=10$ và chiều
cao $10$, nên
\[ S_{FGH}=\frac{10\times 10}{2}=50,\qquad \text{một nửa là } 25. \]
Ta đếm, lần lượt cho từng loại tam giác $GPH$, $FPG$, $HPF$, những vị trí của $P$ cho diện tích $25$.

\begin{recallbox}Hai tam giác có \emph{chung đáy} thì tỉ số diện tích bằng tỉ số hai chiều cao.
Các điểm cách một đường thẳng cùng một khoảng (về cùng một phía) nằm trên một đường thẳng song song
với đường thẳng đó.\end{recallbox}

\step{1}{Tam giác $GPH$}
Tam giác $GPH$ có chung đáy $GH$ với tam giác $FGH$. Muốn diện tích bằng một nửa thì chiều cao hạ từ $P$
phải bằng một nửa chiều cao $10$, tức là $5$. Vậy $P$ nằm trên đường kẻ ngang thứ $5$ (tính từ $GH$).
Trên đường này, cạnh $FG$ đi qua điểm có hoành độ $2{,}5$ và cạnh $FH$ đi qua điểm có hoành độ $7{,}5$
(trung điểm của hai cạnh). Các giao điểm lưới nằm \emph{bên trong} có hoành độ $3,4,5,6,7$:
\[ (3;5),\ (4;5),\ (5;5),\ (6;5),\ (7;5)\quad\Rightarrow\quad 5 \text{ tam giác.} \]

\step{2}{Tam giác $FPG$}
Lần này đáy chung là $FG$. Gọi $M(5;0)$ là trung điểm của $GH$ và $N(7{,}5;5)$ là trung điểm của $FH$.
\begin{proofbox}Tam giác $FMG$ và $FHG$ có chung chiều cao hạ từ $F$, còn đáy $GM=\tfrac12 GH$, nên
$S_{FMG}=\tfrac12 S_{FHG}=25$. Như vậy điểm $M$ cách $FG$ một khoảng đúng bằng nửa khoảng cách từ $H$
đến $FG$. Tương tự với $N$ (trung điểm của $FH$). Do đó $S_{FPG}=25$ khi và chỉ khi $P$ cách $FG$ đúng bằng
$M$, tức là $P$ nằm trên đường thẳng $MN$ (song song với $FG$).\end{proofbox}
Đi từ $G$ đến $F$ ta sang phải $5$, lên trên $10$, tức là ``sang phải $1$, lên $2$''. Đường $MN$ song song với
$FG$ nên từ $M(5;0)$ ta cũng đi ``sang phải $1$, lên $2$'': gặp $(6;2)$, $(7;4)$, rồi tới $N(7{,}5;5)$ (không
phải giao điểm lưới). Hai điểm $(6;2)$, $(7;4)$ nằm bên trong tam giác $\Rightarrow$ $2$ tam giác.

\step{3}{Tam giác $HPF$}
Hình vẽ đối xứng qua đường thẳng đứng $x=5$ nên kết quả như bước 2: $P=(4;2)$ hoặc $(3;4)$
$\Rightarrow$ $2$ tam giác.

\step{4}{Cộng lại}
$5+2+2=9$ tam giác.

\verify{Kiểm tra $P(6;2)$ bằng cách ``bao hình chữ nhật'': hình chữ nhật từ $x=0$ đến $6$, $y=0$ đến $10$ có
diện tích $60$. Bỏ đi ba tam giác vuông ở ngoài tam giác $FPG$: $\tfrac{5\cdot 10}{2}=25$, $\tfrac{1\cdot 8}{2}=4$,
$\tfrac{6\cdot 2}{2}=6$. Còn lại $60-35=25$. Đúng bằng nửa của $50$.}
\begin{warnbox}Đừng quên hai loại tam giác $FPG$ và $HPF$! Nếu chỉ xét tam giác có đáy nằm ngang $GH$
thì sẽ chọn nhầm phương án (A) $5$.\end{warnbox}
Phương án \textbf{(E)}.
\ans{(E) $9$}
}

% ---------------------------------------------------------------- 24
\begin{probbox}{24}{\stars{1}}
Các điểm $T,U,V,W,X,Y$ nằm trên các cạnh của hình vuông $PQRS$ như hình vẽ, với
\[PT=TU=UQ=QV=VW=WR=XS=SY.\]
Phần tô màu chiếm bao nhiêu phần diện tích hình vuông $PQRS$?
\begin{figC}\begin{scope}[x=1.25cm,y=1.25cm]
\fill[tmBlue!35] (1,3)--(2,3)--(1.5,2.5)--(2.5,1.5)--(3,2)--(3,1)--(1,1)--cycle;
\draw[Ink,line width=1pt] (0,0) rectangle (3,3);
\draw[Ink,line width=0.9pt] (1,3)--(1,0) (0,1)--(2,3) (1,3)--(3,1) (1,0)--(3,2) (0,1)--(3,1);
\foreach \x in {0.5,1.5,2.5}{\draw[Ink] (\x-0.03,2.92)--(\x-0.03,3.08) (\x+0.03,2.92)--(\x+0.03,3.08);}
\foreach \y in {0.5,1.5,2.5}{\draw[Ink] (2.92,\y-0.03)--(3.08,\y-0.03) (2.92,\y+0.03)--(3.08,\y+0.03);}
\draw[Ink] (0.47,-0.08)--(0.47,0.08) (0.53,-0.08)--(0.53,0.08);
\draw[Ink] (-0.08,0.47)--(0.08,0.47) (-0.08,0.53)--(0.08,0.53);
\node[above left] at (0,3){$P$};\node[above right] at (3,3){$Q$};\node[below right] at (3,0){$R$};
\node[below left] at (0,0){$S$};\node[above] at (1,3){$T$};\node[above] at (2,3){$U$};
\node[right] at (3,2){$V$};\node[right] at (3,1){$W$};\node[below] at (1,0){$X$};\node[left] at (0,1){$Y$};
\end{scope}\end{figC}
\par\smallskip\noindent
\begin{tabularx}{\linewidth}{@{}XXXXX@{}}(A)\ $\dfrac{5}{18}$&(B)\ $\dfrac13$&(C)\ $\dfrac29$&(D)\ $\dfrac14$&(E)\ $\dfrac16$\end{tabularx}
\end{probbox}
\probsol{
\textbf{Phân tích đề.} Mỗi cạnh hình vuông được chia thành $3$ đoạn bằng nhau. Ta chọn mỗi đoạn dài $1$ đơn
vị, khi đó hình vuông có cạnh $3$ và diện tích $3\times 3=9$. Ý hay: kẻ thêm các đường chia hình vuông thành
lưới $3\times 3$ ô vuông đơn vị (nét chấm), mọi đường trong hình sẽ hiện ra rất ``ngay ngắn''.
\begin{figC}\begin{scope}[x=1.25cm,y=1.25cm]
\fill[tmBlue!12] (1,3)--(3,3)--(3,1)--(1,1)--cycle;
\fill[tmBlue!40] (1,3)--(3,1)--(1,1)--cycle;
\fill[tmOrange!50] (1,3)--(2,3)--(1.5,2.5)--cycle;
\fill[tmGreen!50] (3,2)--(3,1)--(2.5,1.5)--cycle;
\draw[Slate,densely dotted] (0,0) grid[step=1] (3,3);
\draw[Ink,line width=1pt] (0,0) rectangle (3,3);
\draw[Ink,line width=0.9pt] (1,3)--(1,0) (0,1)--(2,3) (1,3)--(3,1) (1,0)--(3,2) (0,1)--(3,1);
\draw[tmOrange!80!black,dashed] (1,2)--(2,3) (1,3)--(2,2);
\draw[tmGreen!70!black,dashed] (2,2)--(3,1) (2,1)--(3,2);
\node[above left] at (0,3){$P$};\node[above right] at (3,3){$Q$};\node[below right] at (3,0){$R$};
\node[below left] at (0,0){$S$};\node[above] at (1,3){$T$};\node[above] at (2,3){$U$};
\node[right] at (3,2){$V$};\node[right] at (3,1){$W$};\node[below] at (1,0){$X$};\node[left] at (0,1){$Y$};
\node[below left] at (1,1){$Z$};\node[font=\footnotesize] at (1.6,1.45){$2$};
\end{scope}\end{figC}

\step{1}{Nhận ra ba mảnh của phần tô màu}
Gọi $Z$ là giao điểm của $TX$ và $YW$. Phần tô màu gồm ba mảnh không chồng lên nhau:
tam giác lớn $TZW$ (xanh đậm), tam giác nhỏ sát cạnh $TU$ (cam) và tam giác nhỏ sát cạnh $VW$ (xanh lá).

\step{2}{Tam giác $TZW$}
$TQWZ$ là hình vuông $2\times 2$ (gồm $4$ ô đơn vị) và $TW$ là đường chéo của nó.
\begin{recallbox}Đường chéo chia hình vuông (hình chữ nhật) thành hai tam giác có diện tích bằng nhau.\end{recallbox}
Vậy $S_{TZW}=\tfrac12\cdot 4=2$.

\step{3}{Hai tam giác nhỏ}
Xét ô đơn vị có cạnh trên là $TU$. Đường $YU$ đi ``sang phải $1$, lên $1$'' qua các góc ô, nên chứa một
đường chéo của ô này; đường $TW$ đi ``sang phải $1$, xuống $1$'' từ $T$, nên chứa đường chéo còn lại.
Hai đường chéo của một hình vuông chia nó thành $4$ tam giác bằng nhau, và tam giác cam là một trong số đó:
diện tích $\tfrac14\cdot 1=\tfrac14$. Hoàn toàn tương tự, tam giác xanh lá nằm trong ô đơn vị có cạnh phải
là $VW$, được cắt bởi hai đường chéo $XV$ và $TW$: diện tích cũng là $\tfrac14$.

\step{4}{Tính phân số}
Diện tích phần tô màu: $2+\tfrac14+\tfrac14=\tfrac52$. Phân số cần tìm:
\[ \frac{5}{2}:9=\frac{5}{2}\cdot\frac{1}{9}=\frac{5}{18}. \]

\begin{warnbox}Đừng ``ước lượng bằng mắt''! Phần tô màu trông gần bằng $\tfrac13$, nhưng tính chính xác
cho $\tfrac{5}{18}$ (nhỏ hơn $\tfrac{6}{18}=\tfrac13$ một chút).\end{warnbox}
\verify{Đặt $S(0;0)$, cạnh $3$; tính diện tích đa giác tô màu bằng cách cộng các tam giác cho cùng kết quả
$\tfrac52$, và $\tfrac{5}{2}=\tfrac{5}{18}\cdot 9$. Đúng.}
Phương án \textbf{(A)}.
\ans{(A) $\dfrac{5}{18}$}
}

% ---------------------------------------------------------------- 25
\begin{probbox}{25}{\relax}
Người ta xếp các hình chữ nhật kích thước $4\times 2$ theo quy tắc: kể từ hình thứ hai, đỉnh trên -- trái
của mỗi hình được đặt đúng vào trung điểm của cạnh dưới của hình ngay phía trên nó (như hình vẽ). Khi xếp
tất cả $10$ hình chữ nhật theo cách này, chu vi của cả hình là bao nhiêu?
\begin{figC}\begin{scope}[x=0.42cm,y=0.42cm]
\foreach \i in {0,1,2,3}{\draw[Ink,line width=0.9pt] (2*\i,-2*\i) rectangle (2*\i+4,-2*\i+2);}
\foreach \i in {0,1,2}{\draw[Ink] (2*\i+1,-2*\i-0.25)--(2*\i+1,-2*\i+0.25) (2*\i+3,-2*\i-0.25)--(2*\i+3,-2*\i+0.25);}
\node[above] at (2,2){$4$};\node[left] at (0,1){$2$};
\foreach \k in {0,1,2}{\fill[Ink] (9+0.6*\k,-6.8-0.6*\k) circle (2.5pt);}
\end{scope}\end{figC}
\choices{$48$}{$64$}{$90$}{$84$}{$100$}
\end{probbox}
\probsol{
\textbf{Phân tích đề.} Mỗi hình mới lùi sang phải $2$ (nửa chiều dài) và xuống dưới $2$ (đúng bằng chiều
rộng). Cả hình có dạng ``cầu thang''. Ta tính chu vi theo hai cách.

\step{1}{Cách 1: đếm từng loại cạnh}
\emph{Cạnh đứng.} Mỗi hình chữ nhật có cạnh trái và cạnh phải dài $2$, và cả hai đều nằm trên đường viền
(không bị hình khác che). Tổng: $10\times(2+2)=40$.

\emph{Cạnh ngang phía trên.} Hình đầu tiên lộ cả cạnh trên: $4$. Mỗi hình sau, nửa trái cạnh trên áp vào
cạnh dưới của hình trước, chỉ lộ nửa phải: $2$. Tổng: $4+9\times 2=22$.

\emph{Cạnh ngang phía dưới.} Mỗi hình trong $9$ hình đầu chỉ lộ nửa trái cạnh dưới: $2$; hình cuối lộ cả
cạnh dưới: $4$. Tổng: $9\times 2+4=22$.

Chu vi: $40+22+22=84$.

\step{2}{Cách 2: ``kéo thẳng'' cầu thang}
\begin{proofbox}Đi dọc đường viền phía trên -- phải từ góc trên -- trái của hình đầu tiên đến góc dưới --
phải của hình cuối, ta chỉ đi sang phải hoặc đi xuống. Các đoạn ngang cộng lại đúng bằng bề ngang cả hình,
các đoạn đứng cộng lại đúng bằng chiều cao cả hình. Đường viền phía dưới -- trái cũng vậy. Do đó chu vi
cầu thang \emph{bằng chu vi hình chữ nhật nhỏ nhất bao quanh nó}.\end{proofbox}
Bề ngang: $4+9\times 2=22$; chiều cao: $10\times 2=20$. Chu vi: $(22+20)\times 2=84$.

\begin{warnbox}Không lấy $10\times$ chu vi một hình ($10\times 12=120$): những đoạn hai hình áp vào nhau
nằm \emph{bên trong}, không thuộc đường viền.\end{warnbox}
\verify{Thử với $2$ hình: cách 1 cho $8+(4+2)+(2+4)=20$; hình bao có kích thước $6\times 4$, chu vi $20$.
Với $n$ hình, chu vi là $2\big((4+2(n-1))+2n\big)=8n+4$; $n=10$ cho $84$. Khớp.}
Phương án \textbf{(D)}.
\ans{(D) $84$}
}

% ---------------------------------------------------------------- 26
\begin{probbox}{26}{\relax}
Một tứ giác có bốn đỉnh là các điểm $A(7;6)$, $B(-5;1)$, $C(-2;-3)$ và $D(10;2)$ trên mặt phẳng tọa độ.
Diện tích tứ giác (tính theo đơn vị diện tích) là
\choices{$60$}{$63$}{$67$}{$70$}{$72$}
\end{probbox}
\probsol{
\textbf{Phân tích đề.} Tứ giác có các cạnh ``xiên'', không có sẵn công thức. Kĩ thuật hay dùng:
\emph{bao tứ giác bằng một hình chữ nhật có cạnh song song với các trục}, rồi trừ đi các tam giác vuông ở
bốn góc.
\begin{figC}\begin{scope}[x=0.4cm,y=0.4cm]
\fill[tmOrange!15] (-5,-3) rectangle (10,6);
\fill[tmBlue!30] (7,6)--(-5,1)--(-2,-3)--(10,2)--cycle;
\draw[Line] (-5,-3) grid[step=1] (10,6);
\draw[->,Slate] (-6,0)--(11,0) node[right]{$x$};\draw[->,Slate] (0,-4)--(0,7) node[above]{$y$};
\draw[tmOrange!80!black,line width=0.9pt] (-5,-3) rectangle (10,6);
\draw[tmBlue!80!black,line width=1.2pt] (7,6)--(-5,1)--(-2,-3)--(10,2)--cycle;
\node[above] at (7,6){$A$};\node[left] at (-5,1){$B$};\node[below] at (-2,-3){$C$};\node[right] at (10,2){$D$};
\node[font=\footnotesize] at (-2.5,4.6){$30$};\node[font=\footnotesize] at (-4.2,-2){$6$};
\node[font=\footnotesize] at (6.5,-1.6){$30$};\node[font=\footnotesize] at (9.1,4.9){$6$};
\end{scope}\end{figC}

\step{1}{Hình chữ nhật bao quanh}
Hoành độ nhỏ nhất là $-5$, lớn nhất là $10$; tung độ nhỏ nhất là $-3$, lớn nhất là $6$.
\begin{recallbox}Khoảng cách giữa hai điểm trên một trục số bằng số lớn trừ số nhỏ: từ $-5$ đến $10$ là
$10-(-5)=15$.\end{recallbox}
Hình chữ nhật có chiều dài $10-(-5)=15$, chiều rộng $6-(-3)=9$, diện tích $15\times 9=135$.

\step{2}{Bốn tam giác vuông ở góc}
\begin{recallbox}Diện tích tam giác vuông $=\tfrac12\times$ tích hai cạnh góc vuông.\end{recallbox}
\begin{itemize}[leftmargin=1.4em,itemsep=1pt]
\item Góc trên -- trái (giữa $B$ và $A$): hai cạnh $7-(-5)=12$ và $6-1=5$: diện tích $\tfrac{12\cdot5}{2}=30$.
\item Góc dưới -- trái (giữa $B$ và $C$): hai cạnh $-2-(-5)=3$ và $1-(-3)=4$: diện tích $\tfrac{3\cdot4}{2}=6$.
\item Góc dưới -- phải (giữa $C$ và $D$): hai cạnh $10-(-2)=12$ và $2-(-3)=5$: diện tích $30$.
\item Góc trên -- phải (giữa $D$ và $A$): hai cạnh $10-7=3$ và $6-2=4$: diện tích $6$.
\end{itemize}
Tổng: $30+6+30+6=72$.

\step{3}{Diện tích tứ giác}
$135-72=63$.

\begin{warnbox}Với tọa độ âm, rất dễ tính nhầm độ dài, chẳng hạn viết $7-5=2$ thay vì $7-(-5)=12$.
Hãy đếm ô trên hình để kiểm tra.\end{warnbox}
\verify{Thực ra $ABCD$ là hình bình hành ($A\to B$ và $D\to C$ đều là ``sang trái $12$, xuống $5$''). Tam giác
$ABD$ có diện tích bằng nửa tứ giác; bao nó bằng hình chữ nhật $15\times 5$ (từ $y=1$ đến $6$):
$75-30-\tfrac{15\cdot 1}{2}-6=31{,}5$. Nhân đôi: $63$. Khớp.}
Phương án \textbf{(B)}.
\ans{(B) $63$}
}

% ---------------------------------------------------------------- 27
\begin{probbox}{27}{\relax}
Hình dưới đây gồm $8$ hình bình hành nhỏ giống hệt nhau ghép lại. Kể cả $8$ hình bình hành nhỏ đó, trong
hình có tất cả bao nhiêu hình bình hành?
\begin{figC}\begin{scope}[x=0.6cm,y=0.6cm]
\foreach \k in {0,...,4}{\draw[Ink,line width=0.9pt] (2.2*\k,0)--(2.2*\k+2,2);}
\foreach \j in {0,1,2}{\draw[Ink,line width=0.9pt] (\j,\j)--(\j+8.8,\j);}
\end{scope}\end{figC}
\choices{$29$}{$30$}{$26$}{$27$}{$28$}
\end{probbox}
\probsol{
\textbf{Phân tích đề.} Mỗi hình bình hành trong hình có hai cạnh nằm trên hai đường \emph{ngang} và hai cạnh
nằm trên hai đường \emph{xiên} của hình. Ta đếm có hệ thống để không sót, không trùng.

\step{1}{Đếm theo chiều cao và chiều ngang (tính bằng số ô)}
\emph{Hình cao $1$ hàng} (có $2$ hàng để chọn). Trong một hàng có $4$ ô liền nhau:
\begin{itemize}[leftmargin=1.4em,itemsep=1pt]
\item rộng $1$ ô: $4$ hình; rộng $2$ ô: $3$ hình; rộng $3$ ô: $2$ hình; rộng $4$ ô: $1$ hình;
\item mỗi hàng có $4+3+2+1=10$ hình, hai hàng có $20$ hình.
\end{itemize}
\emph{Hình cao $2$ hàng}: tương tự có $4+3+2+1=10$ hình.

Tổng: $20+10=30$.

\step{2}{Cách đếm nhanh: chọn đường}
\begin{proofbox}Một hình bình hành được xác định hoàn toàn khi ta chọn $2$ trong $3$ đường ngang (cạnh trên,
cạnh dưới) và $2$ trong $5$ đường xiên (cạnh trái, cạnh phải). Ngược lại, mỗi cách chọn như vậy cho đúng một
hình bình hành.\end{proofbox}
\begin{recallbox}Số cách chọn $2$ đường từ $k$ đường là $\dfrac{k\times(k-1)}{2}$ (mỗi cặp được đếm hai lần
nếu ta chọn ``đường thứ nhất, đường thứ hai'').\end{recallbox}
Chọn $2$ trong $3$ đường ngang: $\tfrac{3\cdot 2}{2}=3$ cách. Chọn $2$ trong $5$ đường xiên:
$\tfrac{5\cdot4}{2}=10$ cách. Số hình bình hành: $3\times 10=30$.

\begin{warnbox}Các bạn hay quên những hình cao $2$ hàng, hoặc quên chính hình lớn nhất (cả hình).\end{warnbox}
\verify{Hai cách đếm cho cùng kết quả $30$.}
Phương án \textbf{(B)}.
\ans{(B) $30$}
}

% ---------------------------------------------------------------- 28
\begin{probbox}{28}{\relax}
Một hình lăng trụ đứng tam giác có thể tích $120$ cm$^3$. Mặt đáy là tam giác vuông có hai cạnh góc vuông
dài $3$ cm và $4$ cm (như hình vẽ). Chiều cao của hình lăng trụ là bao nhiêu xăng-ti-mét?
\begin{figC}\begin{scope}[x=0.75cm,y=0.75cm]
\coordinate (A) at (0,2.8);\coordinate (B) at (1.9,3.8);\coordinate (C) at (4.3,2.2);
\coordinate (A1) at (0,0);\coordinate (B1) at (1.9,1.0);\coordinate (C1) at (4.3,-0.6);
\fill[tmBlue!15] (A)--(B)--(C)--cycle;
\draw[Ink,line width=1pt] (A)--(B)--(C)--cycle (A)--(A1)--(C1)--(C);
\draw[Ink,dashed,line width=0.8pt] (B)--(B1) (A1)--(B1)--(C1);
\draw[Ink] ($(B)!0.3cm!(A)$)--($(B)!0.3cm!(A)+(B)!0.3cm!(C)-(B)$)--($(B)!0.3cm!(C)$);
\draw[Ink] ($(B1)!0.3cm!(A1)$)--($(B1)!0.3cm!(A1)+(B1)!0.3cm!(C1)-(B1)$)--($(B1)!0.3cm!(C1)$);
\node[above left] at ($(A)!0.5!(B)$){$3$ cm};\node[above right] at ($(B)!0.5!(C)$){$4$ cm};
\end{scope}\end{figC}
\choices{$12$}{$20$}{$10$}{$16$}{$8$}
\end{probbox}
\probsol{
\textbf{Phân tích đề.} Biết thể tích và biết đáy (tam giác vuông). Cần tìm chiều cao lăng trụ, tức là độ dài
cạnh bên (nối đáy trên với đáy dưới).

\begin{recallbox}Thể tích hình lăng trụ đứng $=$ diện tích đáy $\times$ chiều cao:
$V=S_{\text{đáy}}\cdot h$. Do đó $h=V:S_{\text{đáy}}$.\end{recallbox}

\step{1}{Diện tích đáy}
Đáy là tam giác vuông, lấy một cạnh góc vuông làm đáy thì cạnh góc vuông kia là chiều cao:
\[ S_{\text{đáy}}=\frac{3\times 4}{2}=6 \text{ (cm}^2). \]

\step{2}{Chiều cao lăng trụ}
$h=120:6=20$ (cm).

\begin{warnbox}Đừng nhầm ``chiều cao của lăng trụ'' (khoảng cách giữa hai đáy) với chiều cao của tam giác đáy.
Cũng đừng quên chia $2$ khi tính diện tích tam giác: nếu lấy $3\times 4=12$ thì sẽ ra $10$ -- phương án (C)
``bẫy''.\end{warnbox}
\verify{$6\times 20=120$ cm$^3$. Đúng.}
Phương án \textbf{(B)}.
\ans{(B) $20$ cm}
}

% ---------------------------------------------------------------- 29
\begin{probbox}{29}{\stars{1}}
Người ta dùng các khối lập phương nhỏ cạnh $1$ để xếp thành một hình hộp chữ nhật kích thước
$8\times 8\times n$ ($n$ là số nguyên dương). Gọi $A$ là diện tích toàn phần của hình hộp, và $B$ là tổng diện
tích toàn phần của tất cả các khối lập phương nhỏ (tính riêng từng khối rồi cộng lại). Tính tổng tất cả các giá
trị của $n$ sao cho $\dfrac{B}{A}$ là một số nguyên.
\choices{$86$}{$90$}{$70$}{$78$}{$96$}
\end{probbox}
\probsol{
\textbf{Phân tích đề.} Ta viết $A$ và $B$ theo $n$, lập tỉ số $\dfrac BA$, rồi tìm những $n$ làm tỉ số
đó nguyên.

\step{1}{Tính $A$}
\begin{recallbox}Hình hộp chữ nhật kích thước $a\times b\times c$ có $6$ mặt, gồm ba cặp mặt bằng nhau, nên
diện tích toàn phần là $2(ab+bc+ca)$.\end{recallbox}
$A=2(8\cdot 8+8\cdot n+8\cdot n)=2(64+16n)=128+32n$.

\step{2}{Tính $B$}
Số khối lập phương nhỏ: $8\cdot 8\cdot n=64n$. Mỗi khối có $6$ mặt, mỗi mặt diện tích $1$, nên
$B=64n\cdot 6=384n$.

\step{3}{Rút gọn tỉ số}
\[ \frac{B}{A}=\frac{384n}{128+32n}=\frac{32\cdot 12n}{32\cdot(4+n)}=\frac{12n}{n+4}. \]
Ta ``tách phần nguyên'': vì $12n=12(n+4)-48$ nên
\[ \frac{12n}{n+4}=\frac{12(n+4)-48}{n+4}=12-\frac{48}{n+4}. \]

\step{4}{Điều kiện để tỉ số là số nguyên}
$12-\dfrac{48}{n+4}$ là số nguyên khi và chỉ khi $\dfrac{48}{n+4}$ là số nguyên, tức là $n+4$ là
\emph{ước} của $48$.
\begin{recallbox}Các ước dương của $48$: $1,2,3,4,6,8,12,16,24,48$.\end{recallbox}
Vì $n\ge 1$ nên $n+4\ge 5$, ta chỉ lấy các ước $6,8,12,16,24,48$, ứng với
\[ n=2,\ 4,\ 8,\ 12,\ 20,\ 44. \]

\step{5}{Tính tổng}
$2+4+8+12+20+44=90$.

\begin{center}\begin{tabular}{|c|c|c|c|c|c|c|}\hline
$n$ & $2$ & $4$ & $8$ & $12$ & $20$ & $44$\\\hline
$B/A=\dfrac{12n}{n+4}$ & $\tfrac{24}{6}=4$ & $\tfrac{48}{8}=6$ & $\tfrac{96}{12}=8$ & $\tfrac{144}{16}=9$ & $\tfrac{240}{24}=10$ & $\tfrac{528}{48}=11$\\\hline
\end{tabular}\end{center}
\begin{warnbox}Đừng quên $n+4$ có thể bằng $48$ (tức $n=44$), một giá trị khá lớn. Và nhớ loại các ước
$1,2,3,4$ vì khi đó $n\le 0$.\end{warnbox}
\verify{Bảng trên cho thấy cả $6$ giá trị đều cho tỉ số nguyên. Thử một giá trị khác, $n=1$:
$\tfrac{12}{5}$ không nguyên. Đúng.}
Phương án \textbf{(B)}.
\ans{(B) $90$}
}

% ---------------------------------------------------------------- 30
\begin{probbox}{30}{\relax}
Hình chữ nhật $PQRS$ được ghép từ sáu hình vuông bằng nhau. Các điểm $U,V,W,X,Y,Z$ là trung điểm của
cạnh các hình vuông như hình vẽ. Tam giác nào dưới đây có diện tích lớn nhất?
\begin{figC}\begin{scope}[x=0.6cm,y=0.6cm]
\draw[Ink,line width=1pt] (0,0) rectangle (6,4);
\draw[Ink,line width=0.9pt] (2,0)--(2,4) (4,0)--(4,4) (0,2)--(6,2);
\foreach \p in {(1,4),(3,4),(4,3),(5,2),(4,1),(6,1)}{\fill[Ink] \p circle (2.2pt);}
\node[above] at (1,4){$V$};\node[above] at (3,4){$W$};\node[right] at (4,3){$X$};
\node[below] at (5,2){$Y$};\node[left] at (4,1){$U$};\node[right] at (6,1){$Z$};
\node[above left] at (0,4){$Q$};\node[above right] at (6,4){$R$};
\node[below left] at (0,0){$P$};\node[below right] at (6,0){$S$};
\end{scope}\end{figC}
\choices{$PVU$}{$PXZ$}{$PVX$}{$PYS$}{$PQW$}
\end{probbox}
\probsol{
\textbf{Phân tích đề.} Ta cần tính (hoặc so sánh) diện tích của $5$ tam giác. Để tránh phân số, chọn cạnh
mỗi hình vuông nhỏ là $2$; khi đó các trung điểm có tọa độ nguyên. Đặt $P(0;0)$, $S(6;0)$, $R(6;4)$, $Q(0;4)$;
suy ra $V(1;4)$, $W(3;4)$, $X(4;3)$, $Y(5;2)$, $U(4;1)$, $Z(6;1)$.

\begin{recallbox}Diện tích tam giác $=\tfrac12\times$ đáy $\times$ chiều cao. Với tam giác ``xiên'', bao nó
bằng hình chữ nhật có cạnh theo lưới rồi trừ các tam giác vuông ở góc (như Bài 26).\end{recallbox}

\step{1}{Hai tam giác có cạnh nằm trên cạnh hình chữ nhật}
\begin{itemize}[leftmargin=1.4em,itemsep=1pt]
\item $PYS$: đáy $PS=6$, chiều cao là khoảng cách từ $Y$ đến $PS$, bằng $2$: $S=\tfrac{6\cdot 2}{2}=6$.
\item $PQW$: đáy $PQ=4$, chiều cao là khoảng cách từ $W$ đến $PQ$, bằng $3$: $S=\tfrac{4\cdot 3}{2}=6$.
\end{itemize}

\step{2}{Ba tam giác ``xiên''}
\begin{itemize}[leftmargin=1.4em,itemsep=1pt]
\item $PVU$: hình chữ nhật bao $4\times 4=16$; trừ $\tfrac{1\cdot 4}{2}=2$, $\tfrac{3\cdot 3}{2}=4{,}5$,
$\tfrac{4\cdot 1}{2}=2$. Còn $16-8{,}5=7{,}5$.
\item $PXZ$: hình chữ nhật bao $6\times 3=18$; trừ $\tfrac{4\cdot3}{2}=6$, $\tfrac{2\cdot2}{2}=2$,
$\tfrac{6\cdot1}{2}=3$. Còn $18-11=7$.
\item $PVX$: hình chữ nhật bao $4\times 4=16$; trừ $\tfrac{1\cdot4}{2}=2$, $\tfrac{3\cdot1}{2}=1{,}5$,
$\tfrac{4\cdot3}{2}=6$. Còn $16-9{,}5=6{,}5$.
\end{itemize}
\begin{figC}\begin{scope}[x=0.6cm,y=0.6cm]
\fill[tmOrange!15] (0,0) rectangle (4,4);
\fill[tmBlue!35] (0,0)--(1,4)--(4,1)--cycle;
\draw[Line] (0,0) grid[step=1] (6,4);
\draw[Ink,line width=1pt] (0,0) rectangle (6,4);
\draw[tmBlue!80!black,line width=1.1pt] (0,0)--(1,4)--(4,1)--cycle;
\node[font=\footnotesize] at (0.35,3.4){$2$};\node[font=\footnotesize] at (3,3){$4{,}5$};\node[font=\footnotesize] at (3,0.35){$2$};
\node[below left] at (0,0){$P$};\node[above] at (1,4){$V$};\node[right] at (4,1){$U$};
\end{scope}\end{figC}

\step{3}{So sánh}
$7{,}5>7>6{,}5>6=6$. Tam giác lớn nhất là $PVU$.

\begin{warnbox}Đừng đoán bằng mắt: tam giác $PXZ$ trông ``to'' vì kéo dài sang tận cạnh $RS$, nhưng diện tích
chỉ là $7$, nhỏ hơn $PVU$.\end{warnbox}
\verify{Nếu lấy cạnh hình vuông là $1$ thì mọi diện tích đều chia cho $4$ ($1{,}875$; $1{,}75$; $1{,}625$;
$1{,}5$; $1{,}5$), thứ tự không đổi.}
Phương án \textbf{(A)}.
\ans{(A) tam giác $PVU$}
}
